% Phase 15 insertion snippets. Current manuscript source is absent.
% Restore approved preamble and citation keys before integration.

% Put this text in the existing abstract environment:
The paper evaluates whether perpetual-futures information improves forecasts of cryptocurrency spot illiquidity beyond spot-market history. In the historical 4-hour Binance benchmark, category-average relative MAE reductions for 24-hour Amihud were 18.56\% for Bitcoin and 13.64\% for Ethereum, but only two of 18 frozen asset--target--horizon tasks passed the registered rule. A technical audit found that unstandardized neural forecasts had extreme absolute errors. A retrospective train-only target-standardization sensitivity materially changed the category averages. A separate development-selected retrospective comparison was inconclusive for Bitcoin and favorable for Ethereum, although the Ethereum system changed from Ridge SpotOnly to ExtraTrees SpotPlusDerivatives and therefore does not isolate the feature-set contribution. Broad same-venue asset, temporal and liquidity-proxy generalization failed its original rule, and the frozen Bybit QML evaluation did not replicate incremental value. The contribution is a transparent empirical evaluation of information-set value, model reliability, selection and transferability, not a new learning algorithm.

\subsection{Historical multi-model benchmark}
The frozen Phase 6 benchmark contains 180 model-level results across 18 asset--target--horizon tasks. Two tasks passed the original registered task rule: BTC and ETH 24-hour Amihud. The equal-category relative MAE reductions were 18.5648\% and 13.6448\%, respectively. These are historical category averages, not comparisons against a single development-selected best SpotOnly system. The complete 180-row benchmark is reported in Supplementary Table S1.

\subsection{Retrospective neural target-scaling sensitivity}
A later audit found that Phase 6 MLP and GRU forecasts were not trained with target standardization and sometimes had errors far above the target scale. Phase 13 retained the selected configurations, features, chronological folds and test origins while standardizing the target using training labels only, inverse-transforming predictions, and scoring on the original $\log(1+Y)$ scale. BTC MLP MAE changed from $0.035047/0.014367$ to $7.809\times10^{-6}/7.976\times10^{-6}$ for SpotOnly/SpotPlusDerivatives; ETH MLP changed from $0.009432/0.007983$ to $1.551\times10^{-5}/1.575\times10^{-5}$. BTC GRU changed from $0.034905/0.015862$ to $6.920\times10^{-6}/7.390\times10^{-6}$; ETH GRU changed from $0.005360/0.000913$ to $1.462\times10^{-5}/1.456\times10^{-5}$. The category averages changed from $18.5648\%$ to $-0.9744\%$ for BTC and from $13.6448\%$ to $0.8081\%$ for ETH. Phase 13 is retrospective and does not amend Phase 6. Phase 6 applied the lower bound $\max(\hat z,0)$ to predictions on the log1p scale; Phase 13 retained its scaled treatment predictions without clipping. This post-processing difference is reported explicitly in Supplementary Table S2.

\subsection{Development-selected strong-baseline reanalysis}
Phase 14 selected a SpotOnly and a SpotPlusDerivatives system separately using pooled chronological development-validation MAE. Define $\Delta_{\mathrm{best}}=\mathrm{MAE}(\mathrm{BEST}_{S0})-\mathrm{MAE}(\mathrm{BEST}_{S1})$, where positive values favor the selected derivatives-eligible system. BTC selected HGB cfg03 for SpotOnly and HGB cfg02 for SpotPlusDerivatives, with MAEs $6.123\times10^{-6}$ and $6.001\times10^{-6}$ ($\Delta_{\mathrm{best}}=1.220\times10^{-7}$; 1.992\%). Its nominal 42-bar interval crosses zero and its two-asset Holm-adjusted $p$-value is 0.067932. ETH selected Ridge cfg03 for SpotOnly and ExtraTrees cfg02 for SpotPlusDerivatives, with MAEs $1.491\times10^{-5}$ and $1.415\times10^{-5}$ ($\Delta_{\mathrm{best}}=7.605\times10^{-7}$; 5.102\%; nominal 42-bar interval $[1.848\times10^{-7},1.175\times10^{-6}]$; Holm-adjusted $p=0.009990$). The ETH result is a favorable retrospective selected-system difference, not a fixed-model feature effect. In saved same-family comparisons, ETH Ridge cfg03 and HGB cfg03 favor S1 descriptively, while selected-config ExtraTrees favors S0; the ExtraTrees configurations differ and no separate inference was registered.

\subsection{Cross-asset, temporal and liquidity-proxy generalization}
The frozen Phase 7--8 broad generalization rule failed. Favorable Amihud category-average point estimates occurred for BTC, ETH and XRP, but only BTC's archived equal-category interval excluded zero in the favorable direction; SOL and BNB point estimates favored SpotOnly. Kyle-style price impact and Corwin--Schultz spread results varied across assets. BTC and ETH category-average differences reversed in 2025 relative to 2024 partial and 2026 partial periods. These are previously inspected same-venue periods and indicate temporal heterogeneity, not independent replication.

\subsection{Historical QML residual corrections and frozen Bybit evaluation}
The selected historical Binance residual configurations favored matched MLP in four of five seeds for BTC Hybrid QNN 4q and ETH QNN 6q, evaluated by complete corrected-forecast MAE through 2023-06-06 04:00 UTC. On the Bybit holdout (1,315 origins per asset), BTC $\Delta_I$ was positive but inconclusive (raw/Holm $p=0.057942$; nominal 42-block interval crosses zero). ETH $\Delta_I$ was negative (raw $p=0.014985$, Holm $p=0.029970$), indicating worse complete-forecast MAE with the QML correction. The adjudicated classification is \texttt{NO\_EXTERNAL\_QML\_REPLICATION}. Registered stronger Bybit classical residual predictions are unavailable, so no QML-versus-strong-classical claim is supported.

% Suggested discussion/limitations/conclusion prose:
The empirical record is sensitive to learner reliability and evaluation target. Large relative neural improvements contributed to favorable historical category averages despite neural absolute errors far above tree and simple baselines. Target scaling materially changed the category statistic, while scaled neural models remained behind the strongest tree references. The Phase 14 ETH result changes the selected model family and does not identify the isolated contribution of derivatives. Mixed asset-, period-, proxy- and venue-level evidence precludes a universal information advantage. All results concern prediction and do not identify causal transmission, execution-cost savings, or trading profit. The QML study used classical statevector simulation and does not establish hardware or computational quantum advantage.
